\subsection{A literal quartic in 209 unrestricted integers}

Write $d=2^{19}$.  The signed-horizon argument permits the source horizon
$t=2+4z$.  The old horizon parameter occurs only in its defining equation,
so replacing $t$ everywhere by $2+4z$ eliminates one coordinate and one
residual.  There remain 209 independently quantified integer coordinates and
149 residuals of degree at most two.

The retained divisibility equation gives
\[
 F=d\,t(2u+1),\qquad \nu_2(t)=1,\qquad \nu_2(F)=20.
\]
Consequently, for every $s=2^k$ with $0\leq k\leq20$, the substitution
$F=sf$ has an integral inverse on every zero of the source system.  Normalize
the growth-divisibility residual by dividing it by $s$.  This is an integral
coefficientwise division after the horizon substitution.  All other
residuals undergo only linear or affine coordinate substitutions, so every
residual remains quadratic.

For complete coordinate precision, number the source coordinates from 0 to
209 and its residuals from 0 to 149.  The source coordinates 0, 2 and 178 are
$t,F,z$, respectively.  Let $I=\{0,\ldots,149\}\setminus\{125\}$ and
define $\Phi_s:\mathbb Z^{209}\longrightarrow\mathbb Z^{210}$ by
\[
 (\Phi_s(\boldsymbol x))_0=2+4x_0,\qquad
 (\Phi_s(\boldsymbol x))_2=sx_2,\qquad
 (\Phi_s(\boldsymbol x))_{178}=x_0.
\]
For $1\leq i\leq177$ with $i\ne2$, its $i$th coordinate is $x_i$; for
$179\leq i\leq209$, it is $x_{i-1}$.  The numbered source rows here refer
to the frozen affine-horizon source.  Put $c_{149,s}=s$ and $c_{j,s}=1$
otherwise.  The sparse principal is the explicit polynomial
\[
 \boxed{\displaystyle
 P_{209}(\boldsymbol x)=
 \sum_{j\in I}\left(\frac{R_j(\Phi_{64}(\boldsymbol x))}
 {c_{j,64}}\right)^2=0,
 \qquad\boldsymbol x\in\mathbb Z^{209}.}
\]
The displayed divisions are divisions of fixed integer coefficient vectors,
not additional predicates or rational-function operations.  The literal
expanded residuals and all collected polynomial coefficients are supplied.

This equation has 5,486 collected monomials and coefficient height $2^{32}$.
The coefficient of $x_0^2x_{208}^2$ is exactly $2^{32}$, proving that degree
four is attained.  Its largest residual coefficient has absolute value
$2^{16}$.  All 209 coordinates occur.

\subsection{A companion with coefficient height $2^{30}$}

Perform the same substitution with $s=128$ and number its 149 residuals
$S_0,\ldots,S_{148}$.  Only two further row operations are needed:
\[
 S_{98}\longmapsto S_{98}+4S_{135},\qquad
 S_{124}\longmapsto S_{124}+6S_{135}.
\]
Leave $S_{135}$ unchanged.  These operations remove the repeated $f^2$
terms from the two indicated rows.  Their inverse subtracts the same
multiples of the unchanged carrier row, so they constitute an integral
determinant-one change of residual basis.  Squaring and summing the resulting
149 rows gives an equivalent equation with 6,654 monomials and coefficient
height $2^{30}$.  Its coefficient of $x_0^2x_{208}^2$ is $2^{30}$, again
certifying exact degree four.

Relative to the incoming 210-coordinate quartic, the sparse equation saves
one coordinate and one residual and reduces height by a factor of 256,
while using 204 additional monomials.  The companion reduces height by a
factor of 1,024 and uses 1,372 additional monomials.  These are explicit
tradeoffs; they do not establish global minimality.  A finite search over
ten power-of-two coordinate scales is retained solely as a reproducible
comparison within the two specified residual bases.

\paragraph{Verification boundary.}
An independent checker, which does not import the producer, verifies every
substitution and row identity and recomputes each collected sum of squares.
All 15 new equation artifacts and both source stages also rebuild byte for
byte.  These checks establish finite algebraic identities and exact metrics.
The signed-horizon lemma and the inherited Pell, binomial, CRT, Horn-flow and
logical-calibration theorems supply the separate unbounded mathematical
proof of the consistency equivalence.
